\documentclass[10pt,a4paper]{amsart}
\usepackage[margin=19mm]{geometry}
\usepackage{amsmath,amssymb,mathtools}
\usepackage[scr=rsfs,cal=euler]{mathalfa}
\usepackage{xfrac}
\usepackage[unicode,hidelinks,bookmarks=false]{hyperref}
\newcommand{\ie}{i.e.\ }
\newcommand{\cf}{cf.\ }
\newcommand{\resp}{respectively\ }
\newcommand{\Subsection}{\subsection}
\providecommand{\nobreakdash}{\nobreak\hbox{-}}
\setlength{\emergencystretch}{2em}
\multlinegap0pt
\usepackage{bm}
\usepackage{bbm}
\usepackage{dsfont}
\usepackage{xspace}
\usepackage{cancel}
\usepackage{mathtools}
\usepackage{amsthm}
\makeatletter
\@ifundefined{theorem}{\newtheorem{theorem}{Theorem}}{}
\makeatother
\DeclareMathOperator*{\argmin}{argmin}
\title[Team Variance Optimization of n-Player Stochastic Games with]{Team Variance Optimization of n-Player Stochastic Games with Separately Controlled Chains}
\author{Li Xia}
\date{Source: arXiv:2507.22335v1 (CC BY 4.0)}
\begin{document}
\maketitle

\noindent\emph{Source and licence.} Team Variance Optimization of n-Player Stochastic Games with Separately Controlled Chains. Version arXiv:2507.22335v1 (CC BY 4.0), distributed under the Creative Commons Attribution 4.0 International licence, as recorded in the arXiv licence metadata for this exact version and shown on the abstract page https://arxiv.org/abs/2507.22335v1. Full rights notice: licenses/arxiv-2507.22335-CC-BY-4.0.txt.\par\smallskip

\noindent\textit{Editorial scope.} Counted unit: the Theorem whose source label is \texttt{None} in the article (source0.tex lines 870--874, source label \texttt{None}), delivered together with the article's own complete original proof of it (source0.tex lines 875--949, one \texttt{proof} environment of 75 lines). No mathematics is written, repaired or re-derived: statement and proof are the article's own text line for line. Required context retained uncounted: eq\_prob0 (lines 436--439); eq\_diffvar (lines 726--731); theorem3 (lines 773--781); theorem4 (lines 791--798); eq\_policyimprov (lines 842--847). Every definition, display and label of those lines is retained unchanged. Omitted from this package and not redistributed: 1--435, 440--725, 732--772, 782--790, 799--841, 848--869, 950--1442. Nothing inside a retained excerpt is omitted or altered: every retained body is delivered exactly as the article's lines are written, so no omission receipt of any kind accompanies this package. Citations: the retained excerpts carry no citation command at all, so no bibliography entry of the article is transcribed and no reference is redistributed. Reference adaptation: none was needed and none was made; every reference inside the delivered unit resolves inside it, so no unresolved reference or citation is produced. Printed slips kept literal: no printed word, symbol, hypothesis or number of the counted statement or of its proof was altered. Layout and class: the article's own class is \texttt{article} with packages including graphicx, color, latexsym, amsfonts, amssymb, bm, bbm, amsthm, amsmath, subfigure, booktabs, boxedminipage, natbib, algorithm; the wrapper is the lane's pinned \texttt{amsart} editorial wrapper at 10pt on A4, which restates the theorem-like environments the retained text needs and adds the article's own macro definitions that the retained text needs (\texttt{\textbackslash argmin}), with extra wrapper packages (bm, bbm, dsfont, xspace, cancel, mathtools). The delivered numbering is the unit's own. Assets: no retained span contains a figure environment, a graphics-inclusion command, a PDF-page inclusion, a table or a TikZ picture; no image file is copied into this package, the package declares no asset dependency and no image file is redistributed with it. Licence: version arXiv:2507.22335v1 of this article is distributed under the Creative Commons Attribution 4.0 International licence, as recorded in the arXiv licence metadata for this exact version and shown on the abstract page https://arxiv.org/abs/2507.22335v1. A full rights notice is delivered at licenses/arxiv-2507.22335-CC-BY-4.0.txt. Characters: every retained excerpt is pure ASCII, so the delivered engine, XeLaTeX with its default Latin Modern text font, reports no missing character.
\begin{flalign}\label{eq_prob0}
&\mathcal M: \hspace{6cm}  J^*_{\sigma} = \min_{u \in \mathcal
U^{\rm HR}} J^{u}_{\sigma} &
\end{flalign}

\begin{equation}\label{eq_diffvar}
J^{u'}_{\sigma} - J^{u}_{\sigma} = \sum_{i=1}^{n} \bm
\pi^{u'_i}_i[(\bm P^{u'_i}_i - \bm P^{u_i}_i)\bm g^{u_i}_i + (\bm
r^{u'_i}_i - \mu^u \bm 1)^2_{\odot} - (\bm r^{u_i}_i - \mu^u \bm
1)^2_{\odot}] - n(\mu^{u'}-\mu^{u})^2.
\end{equation}

\begin{theorem}\label{theorem3}
If we find a new policy $u'_i \in \mathcal U^{\rm SD}_i$ that
ensures the vector $(\bm P^{u'_i}_i - \bm P^{u_i}_i)\bm g^{u_i}_i +
(\bm r^{u'_i}_i - \mu^u \bm 1)^2_{\odot} - (\bm r^{u_i}_i - \mu^u
\bm 1)^2_{\odot} \leq \bm 0$ elementwisely, then we have
$J^{u'}_{\sigma} \leq J^{u}_{\sigma}$. If the inequality holds
strictly in at least one element, then $J^{u'}_{\sigma} <
J^{u}_{\sigma}$.
\end{theorem}

\begin{theorem}\label{theorem4}
The optimal policy $u^*=(u^*_1,u^*_2,\dots,u^*_n)$ of the team
variance minimization problem (\ref{eq_prob0}) must satisfy $\bm
P^{u'_i}_i \bm g^{u^*_i}_i + (\bm r^{u'_i}_i - \mu^{u^*} \bm
1)^2_{\odot} \geq \bm P^{u^*_i}_i \bm g^{u^*_i}_i + (\bm r^{u^*_i}_i
- \mu^{u^*} \bm 1)^2_{\odot}$ elementwisely, for any $u'_i \in
\mathcal U^{\rm SD}_i$ and $i \in \mathcal I$.
\end{theorem}

\begin{equation}\label{eq_policyimprov}
u^{(l+1)}_i(s) = \argmin\limits_{a \in \mathcal A_i}\left\{
(r_i(s,a) - \mu^{u^{(l)}})^2 + \sum_{s' \in \mathcal
S_i}p_i(s'|s,a)g^{u_i^{(l)}}_{i}(s') \right\},  \quad s \in \mathcal
S_i.
\end{equation}

\begin{theorem}
Algorithm~1 will converge to a strictly local minimum or a
first-order stationary point in the space of mixed policies after a
finite number of iterations.
\end{theorem}

\begin{proof}
First, we prove the convergence of Algorithm~1 within a finite
number of iterations. From (\ref{eq_policyimprov}), we can see that
$u_i^{(l+1)}$ and $u_i^{(l)}$ must satisfy the strict inequality in
Theorem~\ref{theorem3} since $u_i^{(l+1)} \neq u_i^{(l)}$ holds for
at least one player $i$. Therefore, we can derive that
$J^{u^{(l+1)}}_{\sigma} < J^{u^{(l)}}_{\sigma}$ and the team
variance is strictly reduced after each iteration of Algorithm~1.
Since the policy space $\mathcal U^{\rm SD}$ is finite, Algorithm~1
will surely stop after a finite number of iterations.

Second, we study the property of the convergence point. We define a
mixed policy $\delta^{u'}_{u}$ that perturbs from the current policy
$u$ to a new policy $u'$ with a mixing probability $\delta$, where
$u,u' \in \mathcal U^{\rm SD}$ and $0\leq \delta \leq 1$. In other
words, if the system adopts the mixed policy $\delta^{u'}_{u}$, each
player~$i \in \mathcal I$ will adopt policy $u_i$ with probability
$1-\delta$ and adopt policy $u'_i$ with probability $\delta$.
Replacing $\delta^{u'}_{u}$ with $u'$ in (\ref{eq_diffvar}) and
taking the differentiating operation through $\delta \rightarrow 0$,
we can derive the \emph{team variance derivative formula} as
follows.
\begin{equation}\label{eq_derivvar}
\frac{\partial J^{\delta^{u'}_{u}}_{\sigma}}{\partial \delta}
\Big|_{\delta=0} = \sum_{i=1}^{n} \bm \pi^{u_i}_i[(\bm P^{u'_i}_i -
\bm P^{u_i}_i)\bm g^{u_i}_i + (\bm r^{u'_i}_i - \mu^u \bm
1)^2_{\odot} - (\bm r^{u_i}_i - \mu^u \bm 1)^2_{\odot}].
\end{equation}
The above derivative can also be denoted as $\nabla J^{u}_{\sigma}$,
which is the gradient of the team variance $J^{u}_{\sigma}$ with
respect to the mixing probability $\delta$ along with a mixing
direction $u \rightarrow u'$ in the policy space $\mathcal U^{\rm
SD}$. When Algorithm~1 stops at a policy, say $u^*$, it indicates
that (\ref{eq_policyimprov}) cannot be improved anymore, i.e.,
\begin{equation}\label{eq29}
\bm P^{u^*_i}_i \bm g^{u^*_i}_i + (\bm r^{u^*_i}_i - \mu^{u^*} \bm
1)^2_{\odot} \leq \bm P^{u'_i}_i \bm g^{u^*_i}_i + (\bm r^{u'_i}_i -
\mu^{u^*} \bm 1)^2_{\odot}, \quad \forall u'_i \in \mathcal U^{\rm
SD}_i, i \in \mathcal I.
\end{equation}
This indicates that the converged point $u^*$ satisfies the
necessary condition of optimal policies, as stated in
Theorem~\ref{theorem4}. Replacing $u$ with $u^*$ in
(\ref{eq_derivvar}) and substituting (\ref{eq29}) into
(\ref{eq_derivvar}), we derive
\begin{equation}\label{eq30}
\frac{\partial J^{\delta^{u'}_{u^*}}_{\sigma}}{\partial \delta}
\Big|_{\delta=0} = \sum_{i=1}^{n} \bm \pi^{u^*_i}_i[(\bm P^{u'_i}_i
- \bm P^{u^*_i}_i)\bm g^{u^*_i}_i + (\bm r^{u'_i}_i - \mu^{u^*} \bm
1)^2_{\odot} - (\bm r^{u^*_i}_i - \mu^{u^*} \bm 1)^2_{\odot}] \geq
0, \quad \forall u' \in \mathcal U^{\rm SD}.
\end{equation}
For more rigorously discussing the sign of the above derivative, we
further partition the policy space $\mathcal U^{\rm SD}$ into two
parts: $\tilde{\mathcal U}^{\rm SD}_{u^*}$ and $\mathcal U^{\rm
SD}\backslash\tilde{\mathcal U}^{\rm SD}_{u^*}$, where
$\tilde{\mathcal U}^{\rm SD}_{u^*}$ is defined as
\begin{equation}\label{eq31}
\tilde{\mathcal U}^{\rm SD}_{u^*} := \{u' :  \frac{\partial
J^{\delta^{u'}_{u^*}}_{\sigma}}{\partial \delta} \Big|_{\delta=0} =
0\}.
\end{equation}
Thus, with (\ref{eq30}) we have
\begin{equation}
\mathcal U^{\rm SD}\backslash\tilde{\mathcal U}^{\rm SD}_{u^*} :=
\{u' :  \frac{\partial J^{\delta^{u'}_{u^*}}_{\sigma}}{\partial
\delta} \Big|_{\delta=0} > 0\}.
\end{equation}
Therefore, we can see that the converged point $u^*$ is a
first-order stationary point in the mixed policy space spanned by
$u' \in \tilde{\mathcal U}^{\rm SD}_{u^*}$ and a strictly local
minimum in the mixed policy space spanned by $u' \in \mathcal U^{\rm
SD}\backslash\tilde{\mathcal U}^{\rm SD}_{u^*}$. Thus, the theorem
is proved.
\end{proof}

\end{document}
